% Generated draft: requires mathematical and source review.
\documentclass[11pt]{article}
\usepackage[margin=25mm]{geometry}
\usepackage{fontspec}
\setmainfont{Latin Modern Roman}
\usepackage{amsmath,amssymb,mathtools,unicode-math}
\setmathfont{Latin Modern Math}
\usepackage{enumitem,needspace,xurl,hyperref}
\setlength{\parindent}{0pt}
\setlength{\parskip}{5pt}
\setlength{\emergencystretch}{3em}
\newcommand{\N}{\mathbb{N}}
\newcommand{\Z}{\mathbb{Z}}
\newcommand{\Q}{\mathbb{Q}}
\newcommand{\R}{\mathbb{R}}
\newcommand{\C}{\mathbb{C}}
\newcommand{\F}{\mathbb{F}}
\newcommand{\A}{\mathbb{A}}
\newcommand{\PP}{\mathbb{P}}
\newcommand{\E}{\mathbb{E}}
\newcommand{\eps}{\varepsilon}
\newcommand{\dd}{\,\mathrm{d}}
\begin{document}
\section*{UnsolvedMath ID: 30006355}
\textbf{Source problem code:} catalogue entry
\subsection{Definitions and mathematical statement}
Let \(V\) be a vertex operator algebra satisfying the source's hypotheses of being strongly finite and \(C_2\)-cofinite. Writing \(Y(u,z)=\sum_{n\in\Z}u_nz^{-n-1}\), define
\[
C_2(V)=\operatorname{span}\{u_{-2}v:u,v\in V\}.
\]
The condition of \(C_2\)-cofiniteness is \(\dim(V/C_2(V))<\infty\), over the ground field of \(V\). The record does not specify that field or define ``strongly finite''; its intended hypotheses must therefore be confirmed before this becomes a fully specified theorem.

Let \(\mathcal C_V\) denote the intended category of \(V\)-modules, and let \(\mathcal B_V\) denote the intended space of genus-one conformal blocks. The precise module category, the permitted insertions, and the definition of this conformal-block space are not supplied. Write \(\mathcal H_V\) for the categorical space on which the source's Hopf-link construction defines an operator \(S_{\mathrm{cat}}\). Write \(S_{\mathrm{mod}}\) for the modular operator on \(\mathcal B_V\) associated with
\[
\tau\longmapsto-\frac{1}{\tau},\qquad
\tau\in\mathbb H=\{z\in\C:\operatorname{Im}z>0\}.
\]
This specifies the transformation of the parameter, but does not specify the action on blocks, including any coordinate, weight, or normalization factors.

The intended comparison problem is to establish, for every \(V\) in the stated class, an appropriate comparison isomorphism
\[
\Phi_V:\mathcal H_V\xrightarrow{\sim}\mathcal B_V
\]
and, with compatible conventions, prove
\[
\Phi_V\circ S_{\mathrm{cat}}
 =S_{\mathrm{mod}}\circ\Phi_V.
\]
Here the existence and meaning of \(\Phi_V\), and the definitions and normalizations of both operators, are unresolved parts of the formulation. The displayed equality expresses the extracted target; it does not assert that these data have been constructed. The record does not settle whether the intended comparison is an equality of linear operators after normalization or an equality only of projective actions.

The retained report passage also asks about the Gainutdinov--Runkel conjecture, an unspecified isomorphism involving torus conformal blocks and the Lyubashenko coend \(L\), and a coend version of sewing-factorization for \(C_2\)-cofinite vertex operator algebras. It does not state that conjecture or the isomorphism precisely. These are contextual questions, whereas the extracted target is the comparison of the two \(S\)-transformations in the strongly finite case. The passage attributes the corresponding strongly rational result to references labelled [Hua08] and [HK10], and refers to the coend through [Lyu96]; full bibliographic details are absent.
\subsection{Short English statement}
For every strongly finite, \(C_2\)-cofinite vertex operator algebra in the source's intended class, identify the categorical torus space with the appropriate genus-one conformal-block space and prove that the Hopf-link \(S\)-operator corresponds to the modular transformation \(\tau\mapsto-1/\tau\). The hypotheses, spaces, comparison map, and normalization conventions require clarification. This is a statement-only draft requiring human mathematical review.
\subsection{Sources}
\begin{itemize}
\item Local UnsolvedMath registry, permanent ID \(30006355\), exact source problem code \texttt{catalogue entry}, title ``Modular S-Transform Equivalence for Strongly Finite Vertex Algebras.'' The supplied record belongs to the Oberwolfach Reports Open Problems collection. No individual proposer is supplied; the proposed year is recorded as 2025.
\item Supplied report citation: \emph{Subfactors and Applications (2026)}, Oberwolfach Reports, pp. 1823--1886, DOI \texttt{10.4171/owr/2025/34}: \url{https://doi.org/10.4171/owr/2025/34}. The citation's 2026 year differs from the 2025 report designation and proposed year; this discrepancy is unresolved.
\item Supplied literature reference: Gui--Zhang, \emph{Sewing-Factorization Theorems}: \url{https://arxiv.org/abs/2503.23995}.
\item Supplied literature reference: Gui--Zhang, \emph{Pseudo-q-traces and coends}: \url{https://arxiv.org/abs/2508.04532}.
\item Supplied literature reference: \emph{Oberwolfach Report 34/2025}: \url{https://ems.press/content/serial-article-files/52239}.
\item Literature assessment dated exactly \texttt{2026-08-22}: the supplied curation labels the record partially solved and reports that Gui and Zhang solved sewing-factorization, essentially identified the Lyubashenko coend interpretation, and proved the Gainutdinov--Runkel conjecture, while the strongly finite comparison of categorical and modular \(S\)-transformations remains open. This is supplied dated curation, not an independently verified literature assessment. Its qualification ``essentially identified'' is retained.
\item The source passage contains the reference labels [Lyu96], [Hua08], and [HK10] without their bibliographic entries. The supplied registry label \texttt{corrected\_verified} is accompanied by a statement that only delimiters, notation, grammar, and typography were repaired; it does not establish mathematical correctness or independently verify the literature claims.
\end{itemize}
\end{document}
